\documentclass[11pt,a4paper]{article}  % Déclare la classe du document.

\newcommand{\chaptermark}[2]{\og #2 \fg{} (#1)} % ligne uniquement pour éviter une erreur en cas d'article
\include{../1ere-preambule}

\usepackage{epic}   % extension et simplification de  l'environnement picture
\usepackage{graphicx}    % permet d'inclure des graphique externe dans un document
\usepackage{arydshln}   %permet d'avoir des lignes en pointillés

\newcommand{\lignes}[1]{
	\setlength{\unitlength}{1mm}
	\vskip 0,4cm\noindent
	\multido{\i=0+1}{#1}{%
		\noindent
		\begin{picture}(150,5)
			\linethickness{0,005mm}
			\put(-5,0){\dashline{1}(175,5)(5,5)}
		\end{picture}\vskip 0,001cm
	}
	\vskip 0cm
}

\rhead{\textcolor{red}{Correction - DS n$^{\circ}01$ - Automatismes - 08/10/2025}}

\newcommand{\va}[1]{\lvert \ #1 \ \rvert}
\newcommand{\vaf}[1]{\left\lvert \ #1\  \right\rvert}

\begin{document}
	\graphicspath{{images/}}
	
%	\noindent Nom: ${\scriptstyle  .........................................................}$ \hskip 1.8cm Prénom: ${\scriptstyle  ...................................................}$ 
	
	\vspace{.5cm}
	
	\begin{center}
		
		
		\fbox{
			\begin{minipage}[]{15cm}
				\begin{center}
					\LARGE{\rule{0.5cm}{0pt}\rule[-0.25cm]{0pt}{0.9cm}
						\textcolor{red}{Correction - DS n$^{\circ}01$ $-$ Automatismes\\
						Proportion et pourcentage} 
					}
				\end{center}
			\end{minipage}
		}
	\end{center}

%\begin{center}\textbf{55min - Barème indicatif}\end{center}
%\begin{center}\textbf{Les élèves avec un tier-temps ne traitent pas les questions avec  le symbole \ding{46}}\end{center}
%\begin{center}\textbf{Calculatrice non autorisée.\\Les résultats devront être justifiés (à l'aide de calcul ou de propriétés).}\end{center}


\vspace*{-0.5cm}
\exerciceDS{(Calculer un pourcentage) $-$ 6 points}

Calculer:
%\begin{multicols}{2}
	\begin{enumMet}
		\item $25\ \%$ de $700$ ;
		\begin{corrBox}
			$25\ \%\times 700=\dfrac{700}{4}=175$		
		\end{corrBox}
		\item \ding{46} $30\ \%$ de $350$;
		\begin{corrBox}
			$30\ \%\times 350=3\times 10\ \% \times 350=3\times 35=105$		
		\end{corrBox}
		\item $0,5\ \%$ de $46\ 000$;
		\begin{corrBox}
			$0,5\ \%\times 46\ 000=\dfrac{1\ \%\times 46\ 000}{2}=\dfrac{460}{2}=230$
		\end{corrBox}
		\item $180\ \%$ de $30$.
		\begin{corrBox}
			$180\ \%\times 30=30\ \%\times 180= 3\times 10\ \%\times 180=3\times 18=54$
		\end{corrBox}
		\item $5\ \%$ de $300$.
		\begin{corrBox}
			$5\ \%\times 300=\dfrac{10\ \%\times 300}{2}=\dfrac{30}{2}=15$
		\end{corrBox}
		\item $14\ \%$ de $600$.
		\begin{corrBox}
			$14\ \%\times 600=10\ \%\times 600+4\times 1\ \%\times 600=60+4\times 6=84$
		\end{corrBox}
	\end{enumMet}
%\end{multicols}

\vspace*{-0.4cm}
\exerciceDS{(Calculer une proportion ou un effectif) $-$ 5 points}
\begin{enumMet}
	\item Calculer $p$ lorsque $n_S = 200$ et $n_E = 800$.\\
	Écrire le résultat sous la forme d'un pourcentage.
	\begin{corrBox}
		$p=\dfrac{n_S}{n_E}=\dfrac{200}{800}=\dfrac{1}{4}=0,25=25\%$		
	\end{corrBox}
	\item \ding{46} Calculer $p$ lorsque $n_S = 20$ et $n_E= 4\ 000$.\\
	Écrire le résultat sous la forme d'un pourcentage.
	\begin{corrBox}
		$p=\dfrac{n_S}{n_E}=\dfrac{20}{4\ 000}=\dfrac{1}{200}=0,005=0,5\%$		
	\end{corrBox}
	\item Calculer $n_S$ lorsque $p = 0,12$ et $n_E= 2\ 000$.
	\begin{corrBox}
		$p=\dfrac{n_S}{n_E} \iff n_S=p\times n_E=2\ 000\times 0,12=2\ 000\times 0,1+2\times 2\ 000\times 0,01=200+2\times 20=240$		
	\end{corrBox}
	\item Calculer $n_E$ lorsque $p = 0,20$ et $n_S = 600$.
	\begin{corrBox}
		$p=\dfrac{n_S}{n_E} \iff n_E=\dfrac{n_S}{p}=\dfrac{600}{0,2}=\dfrac{6\ 000}{2}=3\ 000$		
	\end{corrBox}
	\item Sachant que $40 \%$ d'une somme $S$ vaut $1\ 200$ \EUR{}, calculer $S$.
	\begin{corrBox}
		$S=\dfrac{100\times 1\ 200}{40}=\dfrac{10\times 1\ 200}{4}=10\times 300=3\ 000$		
	\end{corrBox}
\end{enumMet}

\exerciceDS{ $-$ Taux d'évolution entre $Q_{1}$ et $Q_{2}$ $-$ 4 points}
Dans chacun des cas suivants , calculer l'un des trois nombres $Q_{1}, Q_{2}$ ou $t$, connaissant les deux autres. De plus donnner le taux d'évolution sous forme de pourcentage et indiquer s'il s'agit d'une baisse ou d'une hausse.
%	\begin{multicols}{2}
		\begin{enumMet}
		\item \ding{46} $Q_{1}=5 ; Q_{2}=4$.
		\begin{corrBox}
			$t=\dfrac{Q_2-Q_1}{Q_1}=\dfrac{4-5}{5}=-0,2=-20\ \%$. Il s'agit d'une baisse.
		\end{corrBox}
		\item $Q_{1}=10 ; Q_{2}=12,5$.
		\begin{corrBox}
			$t=\dfrac{Q_2-Q_1}{Q_1}=\dfrac{12,5-10}{10}=0,25=25\ \%$. Il s'agit d'une hausse.
		\end{corrBox}
		\item $Q_{1}=20 ; t=-0,10$.
		\begin{corrBox}
			$t=\dfrac{Q_2-Q_1}{Q_1} \iff Q_2=(1+t)Q_1=20\times 0,9=18$. Il s'agit d'une baisse.
		\end{corrBox}
		\item $Q_{2}=140 ; t=0,4$.
		\begin{corrBox}
			$t=\dfrac{Q_2-Q_1}{Q_1} \iff Q_1=\dfrac{Q_2}{1+t}=\dfrac{140}{1,4}=100$. Il s'agit d'une hausse.
		\end{corrBox}
	\end{enumMet}
%	\end{multicols}

\exerciceDS{$-$ 2 points}

Compléter les phrases suivantes.
\begin{multicols}{2}
	\begin{enumMet}
	\item Augmenter de $30 \%$, c'est multiplier par \phantom{$\dfrac{1}{1}$}$\mRouge{1,3}$
	\item Diminuer de $60 \%$, c'est multiplier par \phantom{$\dfrac{1}{1}$}$\mRouge{0,4}$
	\item Multiplier par 4 c'est augmenter de \phantom{$\dfrac{1}{1}$}$\mRouge{300}$\ \%.
	\item Multiplier par 0,35 c'est diminuer de \phantom{$\dfrac{1}{1}$}$\mRouge{65}$\ \%.	
\end{enumMet}
\end{multicols}

\exerciceDS{ $-$ 4,5 points}

\begin{enumMet}
	\item Une personne paie, pour un groupe, une note de restaurant qui s'élève à 180 \EUR{}, avec le service compris de $20 \%$. Quel est le prix des repas sans le service?
	\begin{corrBox}
		$t=20\ \%=0,2\longrightarrow c=1+0,2=1,2$ d'ou $prix=\dfrac{180}{1,2}=150$ \EUR{}		
	\end{corrBox}
	\item \ding{46} Un commerçant calcule ses prix de vente en prenant un bénéfice de $30 \%$ sur ses prix d'achat. Quel est le prix d'achat d'un article qu'il a vendu $260$\EUR{}.
	\begin{corrBox}
		$t=30\ \%=0,3\longrightarrow c=1+0,3=1,3$ d'ou $prix=\dfrac{260}{1,3}=200$ \EUR{}		
	\end{corrBox}
	\item Le prix d'un article soldé est de $66$ \EUR{}. L'étiquette indique \guillemets{$-40 \%$}. Calculer le prix de l'article avant les soldes.
	\begin{corrBox}
		$t=-40\ \%=-0,4\longrightarrow c=1-0,4=0,6$ d'ou $prix=\dfrac{66}{0,6}=110$ \EUR{}		
	\end{corrBox}
\end{enumMet}


\end{document}

